\chapter{Introduction}

\section{Background and context}

Digital fashion recommendation covers several related but distinct tasks,
including the retrieval of visually similar products, the identification of
substitutes, individual-item recommendation and the assessment of whether a
set of garments forms a coherent outfit. Visual resemblance alone is
insufficient for the final task. Two shirts may be substitutes because they
serve a similar role, whereas a shirt, trousers and shoes may be complements
because their different roles form a plausible ensemble. Early image-based
recommendation research formalised this distinction between visual similarity,
substitution and complementarity \parencite{mcauley2015styles}. A modern review
similarly characterises fashion recommendation as a domain in which visual,
semantic, contextual and user-related signals interact, while identifying
evaluation, explanation and bias as continuing challenges
\parencite{deldjoo2023review}.

Outfit compatibility research has moved beyond manually specified
colour or category rules towards learned relationships among items. A
bidirectional long short-term memory model represented outfits as ordered
garment sequences and learned their compatibility \parencite{han2017bilstm}.
Type-aware embedding spaces subsequently distinguished similarity from
compatibility across different item relationships
\parencite{vasileva2018typeaware}. OutfitTransformer then used
task-specific tokens and self-attention to model outfit-level relationships
for compatibility prediction and complementary-item retrieval
\parencite{sarkar2023outfittransformer}. These approaches demonstrate that
complete-outfit recommendation is relational: a candidate should be assessed
in the context of the other garments rather than scored as an isolated item.
They also demonstrate that architecture, category representation, negative
construction and evaluation task affect what a reported score means.

Large-scale vision--language pre-training provides a potential representation
layer for smaller downstream studies. SigLIP uses independent sigmoid losses
over image--text pairs instead of requiring softmax normalisation across a
global batch, and was evaluated as a transferable vision--language
representation by \textcite{zhai2023siglip}. A frozen SigLIP image
encoder can supply reusable features while a smaller model learns
the project-specific relationship among garments. Freezing does not make the
representation neutral or guarantee transfer to fashion data; it instead
reduces the number of trainable parameters and makes controlled comparisons
feasible within the resources of an MSc project.

Throughout this dissertation, \emph{lightweight} describes the project-specific
compatibility head and the cost of downstream training on precomputed features.
It does not describe the size or original pretraining cost of the SigLIP
encoder on which those features depend.

\section{Problem definition and motivation}

This dissertation studies a practical complete-outfit task. A user supplies
one garment, and the system recommends abstract garment types and colours for
the remaining supported positions. The prototype, named \emph{Cove}, uses four
logical slots: an inner top, an optional outer layer, a bottom and shoes. The
selected garment remains fixed. Weather conditions, selected style and the
user-selected clothing range constrain the candidate space, while alternatives
allow the user to retain control over the final choice.

The learned problem is complete-outfit compatibility scoring followed by
candidate ranking. This is separated from candidate feasibility. For example,
weather logic may suppress an outer layer in hot conditions, and slot rules
prevent two bottoms from occupying the same outfit. Such constraints should
not be inferred indirectly from a compatibility score. The deployed system
combines a learned relational score with explicit constraints and a
deterministic fallback.

The candidate representation is also deliberately product-independent.
Recommendations are expressed as abstract garment states defined by slot,
broad type, colour, style and clothing range. The states are derived from
secondary Polyvore training data. This design reflects the intended
interaction: a recommendation can
guide the use of an existing wardrobe or a later purchase without claiming
that one retailer contains the only suitable solution. Product-independent
candidate generation also creates a search problem because compatibility is a
non-additive score over complete combinations. The study consequently examines
whether bounded exact search produces materially different results from an
earlier quota-limited retrieval path.

The motivation is both methodological and practical. A large end-to-end model
would be difficult to train, compare and audit with the available time and
computing resources. A rules-only system is transparent but cannot learn the
relationships observed in outfits. The adopted middle position retains a
frozen visual representation, trains a comparatively small compatibility head
and evaluates the learned component separately from the contextual rules and
search procedure.

Figure~\ref{fig:functional-scope} summarises the system boundary from user and
context inputs to the functions that produce and preserve recommendations. It
also distinguishes learned compatibility ranking from recognition, explicit
constraints and wardrobe management. Cove is treated as a
context-aware complete-outfit recommender system. Search is the mechanism used
to explore and rank feasible candidate combinations, rather than a separate
product-search claim.

\begin{figure}[htbp]
  \centering
  \resizebox{\textwidth}{!}{%
  \begin{tikzpicture}[
    font=\small,
    box/.style={draw=wmgdarkblue!55, rounded corners=2pt, minimum width=43mm,
      minimum height=10mm, align=center, fill=wmgblue!5},
    service/.style={box, fill=wmgblue!13, draw=wmgblue!75, font=\small\bfseries},
    output/.style={box, fill=green!7, draw=green!45!black},
    arrow/.style={-{Latex[length=2.2mm]}, thick, draw=wmgdarkblue!70}
  ]
    \node[font=\bfseries\color{wmgdarkblue}] (inputtitle) {User and context inputs};
    \node[box, below=4mm of inputtitle] (garment) {Selected garment\\image or wardrobe item};
    \node[box, below=2.5mm of garment] (context) {City weather, style\\and clothing range};
    \node[box, below=2.5mm of context] (wardrobe) {Local wardrobe, saved\\pieces and outfits};

    \node[font=\bfseries\color{wmgdarkblue}, right=17mm of inputtitle] (servicetitle) {System functions};
    \node[service, below=4mm of servicetitle] (recognition) {Garment recognition\\and feature extraction};
    \node[service, below=2.5mm of recognition] (generation) {Candidate generation\\and constraint filtering};
    \node[service, below=2.5mm of generation] (ranking) {Complete-outfit\\compatibility ranking};
    \node[service, below=2.5mm of ranking] (management) {Wardrobe storage\\and outfit editing};

    \node[font=\bfseries\color{wmgdarkblue}, right=17mm of servicetitle] (outputtitle) {User-facing outputs};
    \node[output, below=4mm of outputtitle] (primary) {Primary recommendation\\and diverse alternative};
    \node[output, below=2.5mm of primary] (replacement) {Model-ranked colour\\alternatives by position};
    \node[output, below=2.5mm of replacement] (weather) {Weather-aware clothing\\cue and masked absence};
    \node[output, below=2.5mm of weather] (saved) {Saved garment or\\four-position outfit};

    \draw[arrow] (garment.east) -- (recognition.west);
    \draw[arrow] (context.east) -- (generation.west);
    \draw[arrow] (wardrobe.east) -- (management.west);
    \draw[arrow] (recognition.east) -- (primary.west);
    \draw[arrow] (generation.east) -- (weather.west);
    \draw[arrow] (ranking.east) -- (replacement.west);
    \draw[arrow] (management.east) -- (saved.west);
    \draw[arrow] (recognition) -- (generation);
    \draw[arrow] (generation) -- (ranking);
  \end{tikzpicture}}
  \caption{Functional scope of the Cove complete-outfit recommendation system}
  \label{fig:functional-scope}
\end{figure}


\section{Research gap}

The literature establishes that recurrent, type-aware and attention-based
models can learn fashion compatibility
\parencite{han2017bilstm,vasileva2018typeaware,sarkar2023outfittransformer}.
It does not follow that the most
complex published architecture is automatically the most appropriate model
for a smaller, processed dataset, a frozen visual representation and a
resource-constrained application. Reported results are coupled to dataset
splits, negative construction, candidate availability and evaluation protocol
\parencite{herlocker2004evaluation,vasileva2018typeaware}.
The broader fashion-recommender literature also cautions that offline accuracy
alone does not resolve questions of explanation, bias or operational use
\parencite{deldjoo2023review}.

The gap addressed here is bounded and empirical rather than a claim
that outfit recommendation itself is novel. First, it is uncertain whether a
small relational head on frozen vision--language features can outperform
non-learned and controlled learned alternatives under one consistent
pipeline. Second, aggregate compatibility classification does not establish
candidate-completion quality, calibration or generalisation; corrected FITB,
ablation, repeated seeds and overfitting evidence are also required. Third, a
model score does not by itself specify how a practical system should explore a
large abstract candidate space under weather and clothing-range constraints.
The project addresses these questions jointly while keeping model evaluation
distinct from claims about human taste or satisfaction.

\section{Research question}

The primary research question is:

\begin{quote}
To what extent can a small project-specific outfit-compatibility head built on frozen
vision--language embeddings, combined with exact constrained search over
abstract garment states, provide accurate, computationally practical and
auditable complete-outfit recommendations using secondary fashion data?
\end{quote}

Three sub-questions operationalise the model and system components:

\begin{description}[style=nextline,leftmargin=1.2cm,labelindent=0cm]
  \item[\textbf{RQ1.}] How does the proposed compatibility model compare with
        non-learned and controlled learned baselines under consistent data and
        evaluation protocols?
  \item[\textbf{RQ2.}] Which architectural and training choices are supported
        by compatibility classification, corrected FITB, calibration,
        ablation and generalisation evidence?
  \item[\textbf{RQ3.}] What completeness, latency and control trade-offs arise
        when quota-limited retrieval is replaced by exact constrained search
        over product-independent abstract garment states?
\end{description}

The wording \emph{to what extent} requires negative as well as positive
evidence. A high AUC is treated as evidence of discrimination, not as the
percentage of recommendations that users will like. Model ranking, uncertainty,
runtime, data coverage, bias and domain shift all contribute to the
answer.

\section{Research objectives}

The research question is addressed through the following objectives:

\begin{enumerate}[label=\textbf{O\arabic*.}]
  \item Critically synthesise research on fashion compatibility,
        vision--language representation, recommendation evaluation and
        responsible design, and define a reproducible secondary-data method
        consistent with the approved ethical scope.
  \item Develop and evaluate a lightweight project-specific complete-outfit
        compatibility head
        using frozen SigLIP embeddings, including controlled baselines and
        evidence from classification, FITB, calibration, ablation,
        generalisation, imbalance and uncertainty analyses.
  \item Design and implement a product-independent recommendation prototype
        that combines exact candidate search, explicit contextual constraints,
        local wardrobe management, an interaction design intended to support
        clear and accessible use, and a
        deterministic fallback, then verify its functionality, coverage and
        runtime.
  \item Critically evaluate the confidence, limitations, practical value and
        future deployment opportunities supported by the combined model and
        system evidence.
\end{enumerate}

\section{Scope and boundaries}

The investigation is confined to secondary datasets; no human participants
were recruited for interviews, questionnaires or experiments. The project
received confirmation dated 12 May 2026 that ethical approval was waived. The
required research-integrity, information-security and WMG student-ethics
training evidence, together with the waiver confirmation, is reproduced in
Appendices A and B. The separately published Zara dataset retained locally is
used for recognition and feature-domain diagnostics, not as the sole
recommendation source.

The study does not claim to learn an objective or universal definition of
fashion compatibility. Demographic attributes are not collected, so
demographic fairness and user satisfaction cannot be measured
\parencite{mehrabi2021bias,selbst2019fairness}. The interface
asks the user to choose a menswear or womenswear candidate range; it does not
infer gender from a person or garment image \parencite{keyes2018misgendering}.
Sparse pure-menswear training data
preclude a claim that both ranges have equal model reliability. Long-term
personalisation from clicks is also outside scope because defensible online
learning would require exposure logging, feedback-loop controls and a separate
evaluation design \parencite{chaney2018confounding}. Local wardrobe data do not
retrain the compatibility model.

The evaluated representation is limited to an inner top, optional outer layer,
bottom and shoes. Dresses, jumpsuits and other one-piece garments would require
an explicit multi-slot occupancy representation and separate training and
evaluation. Accessories and culturally specific garment structures are not
modelled comprehensively. Official FITB reporting is restricted to questions
representable by these four slots and available frozen embeddings; coverage is
reported so that the filtered result is not presented as a full official
benchmark score.

The work evaluates offline compatibility and system behaviour, not commercial
conversion or human preference. Earlier engineering models use different
candidate or data protocols and are not treated as causal evidence that one
development stage improves another. Controlled claims are based on
same-split baselines, architecture comparisons, structural ablations and the
exact-versus-quota search experiment.

\section{Contributions}

Within these boundaries, the dissertation makes four contributions. First, it
provides a reproducible compatibility pipeline with lightweight downstream
training that separates a
frozen visual representation, a learned relational head and explicit
contextual constraints. Second, it provides a controlled evaluation comprising
non-learned and learned baselines, three seeds, corrected set-identifier-based
FITB labels, structural ablation, calibration and an explicit overfitting
audit. Third, it develops an abstract candidate representation and evaluates
bounded exact search against a quota-limited retrieval path using the same
deployed objective. Fourth, it demonstrates a local-first prototype in which
the user retains the selected garment, chooses the clothing range, receives
weather-aware alternatives and can fall back to deterministic rules if model
artefacts are unavailable.

The contributions are methodological, empirical and artefactual. The comparison
models are controlled adaptations rather than exact reproductions of the
published architectures.

\section{Dissertation structure}

Chapter 2 critically reviews recommender-system foundations, fashion
compatibility, vision--language representations, candidate search, evaluation
and responsible-system concerns. Chapter 3 specifies the secondary-data
pipeline, model, baselines, experimental protocol, metrics, candidate-search system
and ethical safeguards. Chapter 4 reports the model and system results without
extended interpretation. Chapter 5 analyses model behaviour, architectural
evidence, overfitting and search trade-offs. Chapter 6 relates the findings to
previous research, answers the research questions and evaluates contribution,
validity, practical implications and limitations. Chapter 7 concludes by
assessing the objectives and identifying proportionate future work.

\section{Summary}

This chapter has framed complete-outfit recommendation as relational
compatibility scoring followed by constrained candidate search. It has defined
a bounded empirical gap, one primary research question, three sub-questions
and four objectives. The scope excludes participant claims, online
personalisation and unsupported garment structures, while retaining model,
search, runtime, bias and reproducibility evidence. Chapter 2 now evaluates
the literature from which the representation, compatibility and assessment
choices are derived.
